\chapter{The Affine Linear Type Theory of a Reflexive Fuzzy Domain}
\label{ch:affine-linear-type-theory-reflexive-fuzzy-domain}

\section{Introduction}
\label{sec:affine-linear-type-theory-introduction}

The preceding chapters constructed the semantic category of fuzzy domains.
A fuzzy domain is a strict algebra for the fuzzy ideal monad
\[
    T_{\mathbb V}X
    =
    \operatorname{Idl}_{\mathbb V}(X),
\]
and fuzzy-domain morphisms are the \(\mathbb V\)-functors preserving
finite-flat weighted colimits.  The category
\[
    \mathbf{FDom}_{\mathbb V}
\]
of fuzzy domains carries a symmetric monoidal closed structure
\[
    A\boxtimes_{\mathrm{ff}}B
    \dashv
    [B,-]_{\mathrm{ff}}.
\]
Its tensor unit is terminal.  Thus the semantic structure is affine:
there are canonical weakening morphisms
\[
    X\longrightarrow I_{\mathrm{ff}},
\]
but there are no canonical contraction morphisms
\[
    X\longrightarrow X\boxtimes_{\mathrm{ff}}X
\]
in general.

This chapter extracts the corresponding type theory.  The natural syntax is
therefore not cartesian, but affine linear.  Contexts are interpreted by tensor
products, linear function types are interpreted by the internal hom, and
weakening is interpreted by the terminality of the tensor unit.

The central semantic object is a linear reflexive fuzzy domain.  This is a
fuzzy domain \(D\) equipped with a projection pair
\[
    [D,D]_{\mathrm{ff}}\triangleleft D.
\]
Equivalently, there are fuzzy-domain morphisms
\[
    \operatorname{encode}:[D,D]_{\mathrm{ff}}\to D,
    \qquad
    \operatorname{decode}:D\to [D,D]_{\mathrm{ff}},
\]
with
\[
    \operatorname{decode}\operatorname{encode}
    \cong
    \id_{[D,D]_{\mathrm{ff}}}
\]
and a counit comparison
\[
    \operatorname{encode}\operatorname{decode}
    \Rightarrow
    \id_D.
\]
Such a domain carries a canonical affine self-application operation
\[
    D\boxtimes_{\mathrm{ff}}D\longrightarrow D.
\]

The proof theory is enriched over the fuzzy cosmos.  The primitive comparison
judgement is not a scalar inequality, but an evidence-indexed judgement
\[
    \Gamma\vdash_R s\preceq t:A,
    \qquad
    R\in\mathbb V.
\]
Semantically this means a morphism in the fuzzy cosmos
\[
    R
    \longrightarrow
    [\llbracket \Gamma\rrbracket,
     \llbracket A\rrbracket]_{\mathrm{ff}}
    (\llbracket s\rrbracket,\llbracket t\rrbracket).
\]
When \(\mathbb V\) is the posetal fuzzy cosmos associated to a fuzzy
truth-value object
\[
    \Lambda=(L,\leq,\wedge,\vee,\otimes,\multimap,\bot,\top),
\]
this judgement becomes the scalar judgement
\[
    \Gamma\vdash_a s\preceq t:A,
    \qquad
    a\in L,
\]
with semantic meaning
\[
    a
    \leq
    [\llbracket \Gamma\rrbracket,
     \llbracket A\rrbracket]_{\mathrm{ff}}
    (\llbracket s\rrbracket,\llbracket t\rrbracket).
\]

The chapter constructs the raw affine-linear syntax as a locally
\(\mathbb V\)-categorical symmetric affine monoidal closed category, then
applies the fuzzy ideal monad locally to obtain its finite-flat syntactic
envelope.  The main results are soundness, the enriched truth lemma,
initiality of the syntactic envelope, and the posetal soundness-completeness
theorem.

\section{The finite-flat affine semantic base}
\label{sec:finite-flat-affine-semantic-base}

We write
\[
    \mathbf{FDom}_{\mathbb V}
\]
for the ordinary category of fuzzy domains and finite-flat-colimit-preserving
\(\mathbb V\)-functors.  Its tensor product, tensor unit, and internal hom are
denoted
\[
    A\boxtimes_{\mathrm{ff}}B,
    \qquad
    I_{\mathrm{ff}},
    \qquad
    [A,B]_{\mathrm{ff}}.
\]
The closed structure gives natural bijections
\[
    \mathbf{FDom}_{\mathbb V}
    (A\boxtimes_{\mathrm{ff}}B,C)
    \cong
    \mathbf{FDom}_{\mathbb V}
    (A,[B,C]_{\mathrm{ff}}).
\]
The corresponding evaluation morphism is written
\[
    \operatorname{ev}_{A,B}:
    [A,B]_{\mathrm{ff}}\boxtimes_{\mathrm{ff}}A
    \longrightarrow B.
\]

We use the locally \(\mathbb V\)-categorical enhancement of
\(\mathbf{FDom}_{\mathbb V}\).  Its objects are fuzzy domains and its hom
\(\mathbb V\)-categories are
\[
    \mathbf{FDom}_{\mathbb V}^{\mathbb V}(A,B)
    =
    [A,B]_{\mathrm{ff}}.
\]
Thus the objects of the hom \(\mathbb V\)-category are finite-flat-continuous
\(\mathbb V\)-functors
\[
    A\to B,
\]
and the hom-objects between such functors are those of the internal hom
\([A,B]_{\mathrm{ff}}\).

Composition is the finite-flat-continuous bimorphism
\[
    [B,C]_{\mathrm{ff}}
    \boxtimes_{\mathrm{ff}}
    [A,B]_{\mathrm{ff}}
    \longrightarrow
    [A,C]_{\mathrm{ff}},
    \qquad
    (g,f)\longmapsto gf.
\]
The identity at \(A\) is the object
\[
    \id_A\in[A,A]_{\mathrm{ff}}.
\]

Because \(I_{\mathrm{ff}}\) is terminal, each fuzzy domain \(X\) has a canonical
weakening morphism
\[
    {!}_X:X\longrightarrow I_{\mathrm{ff}}.
\]
No diagonal
\[
    X\longrightarrow X\boxtimes_{\mathrm{ff}}X
\]
is part of this structure.

\section{Linear reflexive fuzzy domains}
\label{sec:linear-reflexive-fuzzy-domains}

\begin{definition}[Linear reflexive fuzzy domain]
\label{def:linear-reflexive-fuzzy-domain-type-theory}
A \emph{linear reflexive fuzzy domain} is a fuzzy domain \(D\) equipped with a
fuzzy projection pair
\[
    [D,D]_{\mathrm{ff}}\triangleleft D.
\]
Thus it consists of fuzzy-domain morphisms
\[
    \operatorname{encode}:[D,D]_{\mathrm{ff}}\to D,
    \qquad
    \operatorname{decode}:D\to [D,D]_{\mathrm{ff}},
\]
together with an adjunction
\[
    \operatorname{encode}\dashv \operatorname{decode}
\]
in the locally ordinary support \(2\)-category of fuzzy domains, whose unit is
invertible.  Equivalently,
\[
    \operatorname{decode}\operatorname{encode}
    \cong
    \id_{[D,D]_{\mathrm{ff}}},
\]
and there is a counit comparison
\[
    \operatorname{encode}\operatorname{decode}
    \Rightarrow
    \id_D
\]
satisfying the triangle identities.
\end{definition}

\begin{definition}[Extensional linear reflexive fuzzy domain]
\label{def:extensional-linear-reflexive-fuzzy-domain-type-theory}
A linear reflexive fuzzy domain is \emph{extensional} if the counit comparison
\[
    \operatorname{encode}\operatorname{decode}
    \Rightarrow
    \id_D
\]
is invertible.  Equivalently,
\[
    \operatorname{encode}\operatorname{decode}
    \cong
    \id_D.
\]
\end{definition}

\begin{definition}[Canonical self-application]
\label{def:canonical-self-application-fuzzy-domain}
Let \(D\) be a linear reflexive fuzzy domain.  The \emph{canonical
self-application} morphism is
\[
    \mathsf{app}_D:
    D\boxtimes_{\mathrm{ff}}D
    \longrightarrow
    D
\]
defined by the composite
\[
    D\boxtimes_{\mathrm{ff}}D
    \xrightarrow{\operatorname{decode}\boxtimes 1_D}
    [D,D]_{\mathrm{ff}}\boxtimes_{\mathrm{ff}}D
    \xrightarrow{\operatorname{ev}_{D,D}}
    D.
\]
\end{definition}

\begin{proposition}[Canonical affine applicative structure]
\label{prop:canonical-affine-applicative-structure-fuzzy-domain}
Every linear reflexive fuzzy domain carries the canonical affine applicative
structure
\[
    \mathsf{app}_D:
    D\boxtimes_{\mathrm{ff}}D
    \longrightarrow
    D.
\]
\end{proposition}

\begin{proof}
The morphism \(\operatorname{decode}\boxtimes 1_D\) is a fuzzy-domain morphism
because \(\boxtimes_{\mathrm{ff}}\) is functorial.  The evaluation map is a
fuzzy-domain morphism by the closed structure of
\(\mathbf{FDom}_{\mathbb V}\).  Therefore their composite is a fuzzy-domain
morphism.
\end{proof}

\section{The cosmos-indexed affine linear language}
\label{sec:cosmos-indexed-affine-linear-language}

\subsection{Types and contexts}
\label{subsec:types-and-contexts-cosmos-indexed}

\begin{definition}[Types]
\label{def:affine-linear-types-cosmos}
The types are generated by
\[
    A,B ::= \mathbf 1 \mid D \mid A\otimes B \mid A\multimap B.
\]
Here \(\mathbf 1\) is the tensor unit type, \(D\) is the distinguished
computational type, \(A\otimes B\) is the tensor type, and
\(A\multimap B\) is the linear function type.
\end{definition}

\begin{definition}[Affine linear contexts]
\label{def:affine-linear-contexts-cosmos}
A context is a finite list
\[
    \Gamma=x_1:A_1,\dots,x_n:A_n.
\]
The empty context is denoted
\[
    \varnothing.
\]
Context concatenation is written
\[
    \Gamma,\Delta.
\]
Contexts are identified up to the canonical associativity, unit, and symmetry
isomorphisms generated by the affine symmetric monoidal structure.
\end{definition}

\subsection{Term judgements}
\label{subsec:term-judgements-cosmos-indexed}

Term judgements have the form
\[
    \Gamma\vdash t:A.
\]
The term constructors are those of affine linear natural deduction:

\[
    \frac{}
         {x:A\vdash x:A},
\]

\[
    \frac{}
         {\Gamma\vdash *:\mathbf 1},
\]

\[
    \frac{\Gamma\vdash s:A
          \qquad
          \Delta\vdash t:B}
         {\Gamma,\Delta\vdash s\otimes t:A\otimes B},
\]

\[
    \frac{\Gamma\vdash p:A\otimes B
          \qquad
          \Delta,x:A,y:B\vdash u:C}
         {\Gamma,\Delta\vdash
          \operatorname{let}\ x\otimes y=p\ \operatorname{in}\ u:C},
\]

\[
    \frac{\Gamma,x:A\vdash t:B}
         {\Gamma\vdash \lambda x.t:A\multimap B},
\]

\[
    \frac{\Gamma\vdash f:A\multimap B
          \qquad
          \Delta\vdash u:A}
         {\Gamma,\Delta\vdash f\,u:B}.
\]

Weakening is admissible:
\[
    \frac{\Gamma\vdash t:A}
         {\Gamma,\Delta\vdash t:A}.
\]
Exchange is admissible by the symmetry of the tensor.

There is no contraction rule.  In particular, one does not have a general rule
\[
    \frac{\Gamma,x:A,x:A\vdash t:B}
         {\Gamma,x:A\vdash t:B}.
\]
Such a rule would require a semantic diagonal
\[
    \llbracket A\rrbracket
    \longrightarrow
    \llbracket A\rrbracket
    \boxtimes_{\mathrm{ff}}
    \llbracket A\rrbracket,
\]
which is not available in the affine fuzzy-domain setting.

The ordinary beta and eta equations for tensor and linear function types are
included:
\[
    \operatorname{let}\ x\otimes y=(s\otimes t)\ \operatorname{in}\ u
    \equiv
    u[s/x,t/y],
\]
\[
    \operatorname{let}\ x\otimes y=p\ \operatorname{in}\ x\otimes y
    \equiv
    p,
\]
\[
    (\lambda x.t)\,u
    \equiv
    t[u/x],
\]
\[
    \lambda x.f\,x
    \equiv
    f,
    \qquad x\notin \operatorname{FV}(f).
\]
The unit type is terminal:
\[
    \Gamma\vdash u:\mathbf 1
    \quad\Longrightarrow\quad
    u\equiv *.
\]

\subsection{Encode, decode, and derived self-application}
\label{subsec:encode-decode-derived-self-application-cosmos}

The distinguished type \(D\) carries primitive term formers
\[
    \frac{\Gamma\vdash f:D\multimap D}
         {\Gamma\vdash \operatorname{encode}(f):D},
\]
and
\[
    \frac{\Gamma\vdash d:D}
         {\Gamma\vdash \operatorname{decode}(d):D\multimap D}.
\]

The reflexive beta law is invertible:
\[
    \operatorname{decode}(\operatorname{encode}(f))
    \equiv
    f
    :
    D\multimap D.
\]
The reflexive eta law is one-sided:
\[
    \operatorname{encode}(\operatorname{decode}(d))
    \preceq
    d
    :
    D.
\]
Its precise enriched meaning is given below as a global comparison cell.

Define derived self-application by
\[
    s\bullet t
    :=
    \operatorname{decode}(s)(t).
\]
Thus from
\[
    \Gamma\vdash s:D,
    \qquad
    \Delta\vdash t:D,
\]
one obtains
\[
    \Gamma,\Delta\vdash s\bullet t:D.
\]

If
\[
    \Gamma,x:D\vdash t:D,
\]
define derived \(D\)-abstraction by
\[
    \lambda^D x.t
    :=
    \operatorname{encode}(\lambda x.t).
\]
Then
\[
    \Gamma\vdash \lambda^D x.t:D.
\]
The derived beta law is
\[
    (\lambda^D x.t)\bullet u
    \equiv
    t[u/x].
\]
The derived eta approximation is
\[
    \lambda^D x.(d\bullet x)
    \preceq
    d,
    \qquad
    x\notin \operatorname{FV}(d).
\]

\section{Cosmos-indexed approximation judgements}
\label{sec:cosmos-indexed-approximation-judgements}

\subsection{Evidence-indexed comparison}
\label{subsec:evidence-indexed-comparison}

\begin{definition}[Cosmos-indexed approximation judgement]
\label{def:cosmos-indexed-approximation-judgement}
Let \(R\in\mathbb V\).  A \emph{cosmos-indexed approximation judgement} has
the form
\[
    \Gamma\vdash_R s\preceq t:A.
\]
It is read as: \(R\) is an object of evidence witnessing that \(s\)
approximates \(t\).
\end{definition}

The intended semantics is a morphism
\[
    R
    \longrightarrow
    [\llbracket \Gamma\rrbracket,
     \llbracket A\rrbracket]_{\mathrm{ff}}
    (\llbracket s\rrbracket,\llbracket t\rrbracket)
\]
in the fuzzy cosmos.

\subsection{Structural rules for evidence}
\label{subsec:structural-rules-evidence}

Reflexivity is witnessed by the tensor unit:
\[
    \frac{}
         {\Gamma\vdash_I t\preceq t:A}.
\]

Transitivity composes evidence by tensor:
\[
    \frac{\Gamma\vdash_R r\preceq s:A
          \qquad
          \Gamma\vdash_S s\preceq t:A}
         {\Gamma\vdash_{S\otimes R} r\preceq t:A}.
\]
The order \(S\otimes R\) matches the enriched composition map
\[
    C(s,t)\otimes C(r,s)
    \longrightarrow
    C(r,t).
\]

Evidence is contravariantly functorial in the evidence object:
\[
    \frac{q:Q\to R
          \qquad
          \Gamma\vdash_R s\preceq t:A}
         {\Gamma\vdash_Q s\preceq t:A}.
\]

Small coproducts aggregate evidence:
\[
    \frac{\Gamma\vdash_{R_i} s\preceq t:A\quad(i\in I)}
         {\Gamma\vdash_{\coprod_{i\in I}R_i} s\preceq t:A}.
\]
The induced comparison is obtained by the universal property of the coproduct.

\subsection{Congruence rules}
\label{subsec:congruence-rules-cosmos}

All term formers are enriched functorial.

Tensor introduction is monotone:
\[
    \frac{\Gamma\vdash_R s\preceq s':A
          \qquad
          \Delta\vdash_S t\preceq t':B}
         {\Gamma,\Delta\vdash_{R\otimes S}
          s\otimes t\preceq s'\otimes t':A\otimes B}.
\]

Linear application is monotone:
\[
    \frac{\Gamma\vdash_R f\preceq g:A\multimap B
          \qquad
          \Delta\vdash_S u\preceq v:A}
         {\Gamma,\Delta\vdash_{R\otimes S}
          f\,u\preceq g\,v:B}.
\]

Linear abstraction is enriched functorial:
\[
    \frac{\Gamma,x:A\vdash_R s\preceq t:B}
         {\Gamma\vdash_R
          \lambda x.s\preceq \lambda x.t:A\multimap B}.
\]

Encode is enriched functorial:
\[
    \frac{\Gamma\vdash_R f\preceq g:D\multimap D}
         {\Gamma\vdash_R
          \operatorname{encode}(f)\preceq\operatorname{encode}(g):D}.
\]

Decode is enriched functorial:
\[
    \frac{\Gamma\vdash_R d\preceq e:D}
         {\Gamma\vdash_R
          \operatorname{decode}(d)\preceq\operatorname{decode}(e):D\multimap D}.
\]

Self-application is enriched functorial:
\[
    \frac{\Gamma\vdash_R s\preceq s':D
          \qquad
          \Delta\vdash_S t\preceq t':D}
         {\Gamma,\Delta\vdash_{R\otimes S}
          s\bullet t\preceq s'\bullet t':D}.
\]

\subsection{Beta and eta comparisons}
\label{subsec:beta-eta-comparisons-cosmos}

The ordinary beta and eta laws for the genuine linear function type hold by
invertible \(I\)-evidence in both directions:
\[
    \Gamma,\Delta
    \vdash_I
    (\lambda x.t)\,u
    \preceq
    t[u/x]
    :
    B,
\]
\[
    \Gamma,\Delta
    \vdash_I
    t[u/x]
    \preceq
    (\lambda x.t)\,u
    :
    B.
\]
Likewise,
\[
    \Gamma
    \vdash_I
    \lambda x.f\,x
    \preceq
    f
    :
    A\multimap B,
\]
and
\[
    \Gamma
    \vdash_I
    f
    \preceq
    \lambda x.f\,x
    :
    A\multimap B,
\]
where \(x\notin\operatorname{FV}(f)\).

The reflexive beta law holds by invertible \(I\)-evidence:
\[
    \Gamma
    \vdash_I
    \operatorname{decode}(\operatorname{encode}(f))
    \preceq
    f
    :
    D\multimap D,
\]
and
\[
    \Gamma
    \vdash_I
    f
    \preceq
    \operatorname{decode}(\operatorname{encode}(f))
    :
    D\multimap D.
\]

The reflexive eta law is one-sided:
\[
    \Gamma
    \vdash_I
    \operatorname{encode}(\operatorname{decode}(d))
    \preceq
    d
    :
    D.
\]
Equivalently,
\[
    \Gamma
    \vdash_I
    \lambda^D x.(d\bullet x)
    \preceq
    d
    :
    D.
\]

\begin{definition}[Definitional equality]
\label{def:definitional-equality-cosmos}
Define
\[
    \Gamma\vdash s\equiv t:A
\]
to mean that both
\[
    \Gamma\vdash_I s\preceq t:A
    \qquad\text{and}\qquad
    \Gamma\vdash_I t\preceq s:A
\]
are derivable, and that the corresponding comparison cells are inverse in the
support groupoid of the relevant hom \(\mathbb V\)-category.
\end{definition}

\section{The raw locally \(\mathbb V\)-categorical syntax}
\label{sec:raw-locally-V-categorical-syntax}

The syntax is organized as a locally \(\mathbb V\)-categorical affine closed
category.  Its objects are contexts.  Its hom \(\mathbb V\)-categories are
categories of affine substitutions.

\subsection{Term \(\mathbb V\)-categories}
\label{subsec:term-V-categories}

For each context \(\Gamma\) and type \(A\), let
\[
    \mathsf{Tm}_D(\Gamma,A)
\]
be the \(\mathbb V\)-category presented as follows.

Its objects are terms
\[
    \Gamma\vdash t:A.
\]
For terms \(s,t\), the hom-object
\[
    \mathsf{Tm}_D(\Gamma,A)(s,t)\in\mathbb V
\]
is generated by the cosmos-indexed approximation judgements
\[
    \Gamma\vdash_R s\preceq t:A.
\]
A derivation of such a judgement is equivalently a morphism
\[
    R\longrightarrow \mathsf{Tm}_D(\Gamma,A)(s,t)
\]
in \(\mathbb V\).

The identity morphism
\[
    I\to \mathsf{Tm}_D(\Gamma,A)(t,t)
\]
is generated by reflexivity.  Composition
\[
    \mathsf{Tm}_D(\Gamma,A)(s,t)
    \otimes
    \mathsf{Tm}_D(\Gamma,A)(r,s)
    \longrightarrow
    \mathsf{Tm}_D(\Gamma,A)(r,t)
\]
is generated by transitivity.

The equations imposed on the presentation are exactly the identity,
associativity, functoriality, beta, eta, and triangle equations of the
approximation calculus.

\begin{lemma}[Term categories are \(\mathbb V\)-categories]
\label{lem:term-categories-are-V-categories}
For every context \(\Gamma\) and type \(A\), the presented structure
\[
    \mathsf{Tm}_D(\Gamma,A)
\]
is a small \(\mathbb V\)-category.
\end{lemma}

\begin{proof}
The reflexivity rule supplies the enriched identities.  The transitivity rule
supplies the enriched composition.  The identity and associativity equations
imposed on derivations are precisely the enriched category axioms.  Smallness
follows because terms and formal derivations are generated by a small grammar
from small collections of variables, types, and rules.
\end{proof}

\subsection{Affine substitutions}
\label{subsec:affine-substitutions}

Let
\[
    \Gamma=x_1:A_1,\dots,x_m:A_m
\]
and
\[
    \Delta=y_1:B_1,\dots,y_n:B_n
\]
be contexts.

An affine substitution
\[
    \sigma:\Gamma\Rightarrow\Delta
\]
assigns to each \(y_i:B_i\) a term
\[
    \Gamma_i\vdash t_i:B_i,
\]
where the contexts \(\Gamma_1,\dots,\Gamma_n\) form an affine decomposition
of \(\Gamma\): each variable of \(\Gamma\) occurs in at most one \(\Gamma_i\),
and variables not appearing in any \(\Gamma_i\) are discarded by weakening.
Exchange is accounted for by the symmetry of the tensor.

For two substitutions
\[
    \sigma=(t_1,\dots,t_n),
    \qquad
    \tau=(u_1,\dots,u_n)
\]
with a common affine source decomposition
\[
    \Gamma_1,\dots,\Gamma_n,
\]
define their comparison object by
\[
    \mathsf{Sub}_D(\Gamma,\Delta)(\sigma,\tau)
    =
    \bigotimes_{i=1}^n
    \mathsf{Tm}_D(\Gamma_i,B_i)(t_i,u_i),
\]
with the empty tensor interpreted as \(I\).  For substitutions represented
using different affine decompositions, the comparison object is obtained as
the coproduct over common affine refinements.  The universal maps from the
refined presentations implement exchange and weakening coherence.

Thus for every pair of contexts \(\Gamma,\Delta\) we obtain a
\(\mathbb V\)-category
\[
    \mathsf{Sub}_D(\Gamma,\Delta).
\]

\begin{lemma}[Substitution homs]
\label{lem:substitution-homs-V-categories}
For every pair of contexts \(\Gamma,\Delta\), the substitutions
\[
    \Gamma\Rightarrow\Delta
\]
with the comparison objects above form a small \(\mathbb V\)-category
\[
    \mathsf{Sub}_D(\Gamma,\Delta).
\]
\end{lemma}

\begin{proof}
Identities are given componentwise by the identities in the term
\(\mathbb V\)-categories.  Composition is given componentwise by composition
in the term \(\mathbb V\)-categories, followed by the coherence maps induced
by affine refinement.  The equations imposed on affine decompositions ensure
independence of the chosen refinement.  Associativity and the unit laws reduce
to the corresponding enriched category axioms for the term categories.
\end{proof}

\subsection{The raw syntactic category}
\label{subsec:raw-syntactic-category}

\begin{definition}[The raw syntactic category]
\label{def:raw-syntactic-category}
Let
\[
    \mathcal L_D
\]
be the locally \(\mathbb V\)-categorical category defined as follows.

Its objects are affine linear contexts.  Its hom \(\mathbb V\)-category is
\[
    \mathcal L_D(\Gamma,\Delta)
    =
    \mathsf{Sub}_D(\Gamma,\Delta).
\]
Composition
\[
    \mathcal L_D(\Delta,\Theta)
    \otimes
    \mathcal L_D(\Gamma,\Delta)
    \longrightarrow
    \mathcal L_D(\Gamma,\Theta)
\]
is substitution.  Identities are identity substitutions.
\end{definition}

\begin{lemma}[Substitution is enriched functorial]
\label{lem:substitution-enriched-functorial}
Composition by substitution defines a \(\mathbb V\)-functor
\[
    \mathcal L_D(\Delta,\Theta)
    \otimes
    \mathcal L_D(\Gamma,\Delta)
    \longrightarrow
    \mathcal L_D(\Gamma,\Theta).
\]
\end{lemma}

\begin{proof}
On objects, composition is ordinary affine simultaneous substitution.
On hom-objects, the substitution and approximation-substitution rules send
componentwise comparison evidence for
\[
    \sigma\preceq\sigma'
    \qquad\text{and}\qquad
    \tau\preceq\tau'
\]
to comparison evidence for
\[
    \tau\sigma\preceq\tau'\sigma'.
\]
The evidence object is the tensor product of the input evidence objects,
exactly as required for a \(\mathbb V\)-functor out of the tensor product of
hom \(\mathbb V\)-categories.  The compatibility equations are precisely the
substitution associativity and identity equations of the syntax.
\end{proof}

\begin{proposition}[The raw syntactic category]
\label{prop:raw-syntactic-category}
The structure
\[
    \mathcal L_D
\]
is a locally \(\mathbb V\)-categorical category.
\end{proposition}

\begin{proof}
The homs are \(\mathbb V\)-categories by
Lemma~\ref{lem:substitution-homs-V-categories}.  Composition is enriched
functorial by Lemma~\ref{lem:substitution-enriched-functorial}.  The identity
and associativity axioms are the ordinary identity and associativity laws for
affine simultaneous substitution, together with their comparison-cell
coherence equations.
\end{proof}

\section{Affine closed structure of the raw syntax}
\label{sec:affine-closed-structure-raw-syntax}

\begin{definition}[Tensor and unit in raw syntax]
\label{def:tensor-unit-raw-syntax}
The tensor product in
\[
    \mathcal L_D
\]
is context concatenation:
\[
    \Gamma\otimes\Delta
    :=
    \Gamma,\Delta.
\]
The tensor unit is the empty context
\[
    \varnothing.
\]
The symmetry, associativity, and unit constraints are generated by exchange
and the structural equations of affine contexts.
\end{definition}

\begin{lemma}[Affinity of the raw syntax]
\label{lem:affinity-raw-syntax}
The tensor unit
\[
    \varnothing
\]
is terminal in the raw syntactic category.
\end{lemma}

\begin{proof}
For every context \(\Gamma\), there is a unique affine substitution
\[
    \Gamma\Rightarrow\varnothing,
\]
namely the empty tuple of terms.  Its comparison object with itself is the
empty tensor, hence \(I\).  Therefore the empty context is terminal in the
support and enriched senses.
\end{proof}

\begin{definition}[Internal hom in raw syntax]
\label{def:internal-hom-raw-syntax}
For types \(A,B\), the internal hom is represented by the type
\[
    A\multimap B.
\]
For contexts, the closed structure is generated by the bijective
correspondence between terms
\[
    \Gamma,x:A\vdash t:B
\]
and terms
\[
    \Gamma\vdash \lambda x.t:A\multimap B.
\]
\end{definition}

\begin{proposition}[Raw affine closed structure]
\label{prop:raw-affine-closed-structure}
The locally \(\mathbb V\)-categorical category
\[
    \mathcal L_D
\]
is symmetric affine monoidal closed.
\end{proposition}

\begin{proof}
The tensor and unit are given by context concatenation and the empty context.
The exchange, associativity, and unit equations give the symmetric monoidal
structure.  Lemma~\ref{lem:affinity-raw-syntax} gives affinity.

For closedness, abstraction and application define mutually inverse enriched
functors between the relevant hom \(\mathbb V\)-categories:
\[
    \mathcal L_D(\Gamma\otimes A,B)
    \cong
    \mathcal L_D(\Gamma,A\multimap B).
\]
The beta and eta laws for the genuine linear function type are precisely the
triangle identities of this enriched tensor--hom adjunction.  Functoriality on
comparison evidence is exactly the abstraction and application congruence
rules.
\end{proof}

\section{The reflexive object in raw syntax}
\label{sec:reflexive-object-raw-syntax}

The distinguished type \(D\) is an object of \(\mathcal L_D\).  The primitive
term formers encode and decode give morphisms
\[
    \operatorname{encode}:D\multimap D\longrightarrow D,
\]
and
\[
    \operatorname{decode}:D\longrightarrow D\multimap D
\]
in the support category of \(\mathcal L_D\).

Using the closed structure, the type \(D\multimap D\) represents the internal
hom from \(D\) to \(D\).  Thus the pair above is the syntactic analogue of
\[
    [D,D]\triangleleft D.
\]

\begin{definition}[Raw syntactic reflexive structure]
\label{def:raw-syntactic-reflexive-structure}
The raw syntactic reflexive structure on \(D\) consists of the support maps
\[
    \operatorname{encode}:D\multimap D\to D,
    \qquad
    \operatorname{decode}:D\to D\multimap D,
\]
together with:

\begin{enumerate}
    \item invertible comparison cells
    \[
        \operatorname{decode}\operatorname{encode}
        \cong
        \id_{D\multimap D};
    \]

    \item a counit comparison cell
    \[
        \operatorname{encode}\operatorname{decode}
        \Rightarrow
        \id_D;
    \]

    \item the triangle equations for the adjunction
    \[
        \operatorname{encode}\dashv\operatorname{decode}.
    \]
\end{enumerate}
\end{definition}

\begin{proposition}[Raw syntactic linear reflexivity]
\label{prop:raw-syntactic-linear-reflexivity}
The distinguished object \(D\) of \(\mathcal L_D\) is a linear reflexive object:
\[
    [D,D]_{\mathcal L_D}\triangleleft D.
\]
\end{proposition}

\begin{proof}
The closed structure identifies
\[
    [D,D]_{\mathcal L_D}
\]
with the type \(D\multimap D\).  The maps
\(\operatorname{encode}\) and \(\operatorname{decode}\) are the primitive
syntactic maps.  The reflexive beta comparison supplies the invertible unit
\[
    \operatorname{decode}\operatorname{encode}
    \cong
    \id_{D\multimap D},
\]
and the reflexive eta comparison supplies the counit
\[
    \operatorname{encode}\operatorname{decode}
    \Rightarrow
    \id_D.
\]
The triangle equations are imposed as equations of comparison cells in the raw
syntactic presentation.  Hence the pair is a projection pair.
\end{proof}

\section{Finite-flat syntactic envelope}
\label{sec:finite-flat-syntactic-envelope}

The raw syntax is locally \(\mathbb V\)-categorical, but its homs are not
required to be fuzzy domains.  We now apply the fuzzy ideal monad locally.

\begin{definition}[Finite-flat syntactic envelope]
\label{def:finite-flat-syntactic-envelope}
The \emph{finite-flat syntactic envelope}
\[
    \mathcal L_D^{\mathrm{ff}}
\]
has the same objects as
\[
    \mathcal L_D.
\]
For contexts \(\Gamma,\Delta\), its hom fuzzy domain is
\[
    \mathcal L_D^{\mathrm{ff}}(\Gamma,\Delta)
    =
    \operatorname{Idl}_{\mathbb V}
    \bigl(\mathcal L_D(\Gamma,\Delta)\bigr).
\]
The Yoneda embedding of each hom \(\mathbb V\)-category is written
\[
    Y_{\Gamma,\Delta}:
    \mathcal L_D(\Gamma,\Delta)
    \longrightarrow
    \mathcal L_D^{\mathrm{ff}}(\Gamma,\Delta).
\]
\end{definition}

Let
\[
    c_{\Gamma,\Delta,\Theta}:
    \mathcal L_D(\Delta,\Theta)\otimes
    \mathcal L_D(\Gamma,\Delta)
    \longrightarrow
    \mathcal L_D(\Gamma,\Theta)
\]
be raw composition.  Given finite-flat weights
\[
    P\in
    \operatorname{Idl}_{\mathbb V}
    \bigl(\mathcal L_D(\Delta,\Theta)\bigr),
\]
and
\[
    Q\in
    \operatorname{Idl}_{\mathbb V}
    \bigl(\mathcal L_D(\Gamma,\Delta)\bigr),
\]
their composite is defined by external product followed by lower direct image:
\[
    P\circ_{\mathrm{ff}}Q
    =
    (c_{\Gamma,\Delta,\Theta})_!
    (P\boxtimes Q).
\]
Thus composition in the finite-flat envelope is
\[
    \mathcal L_D^{\mathrm{ff}}(\Delta,\Theta)
    \boxtimes_{\mathrm{ff}}
    \mathcal L_D^{\mathrm{ff}}(\Gamma,\Delta)
    \longrightarrow
    \mathcal L_D^{\mathrm{ff}}(\Gamma,\Theta).
\]

\begin{lemma}[Finite-flat composition]
\label{lem:finite-flat-composition-syntactic-envelope}
Composition in
\[
    \mathcal L_D^{\mathrm{ff}}
\]
is finite-flat-continuous in each variable.
\end{lemma}

\begin{proof}
External products of finite-flat weights are finite-flat.  Lower direct image
along a \(\mathbb V\)-functor preserves finite-flat weights.  Hence the
displayed composite is again a finite-flat weight.

Finite-flat continuity in each variable follows from the Fubini law for
finite-flat weights and functoriality of lower direct image:
iterated finite-flat colimits of external products are computed by the
corresponding finite-flat colimit of composites.
\end{proof}

\begin{lemma}[Associativity and units in the envelope]
\label{lem:associativity-units-finite-flat-envelope}
The finite-flat composition above is associative and unital.
\end{lemma}

\begin{proof}
Associativity of raw substitution gives equality of the two raw composition
functors
\[
    \mathcal L_D(\Theta,\Xi)\otimes
    \mathcal L_D(\Delta,\Theta)\otimes
    \mathcal L_D(\Gamma,\Delta)
    \longrightarrow
    \mathcal L_D(\Gamma,\Xi).
\]
Taking external products of finite-flat weights and then lower direct image
along these equal raw composites gives equal results.

The unit laws follow similarly from the raw identity substitution.  The
principal ideal of the identity substitution acts as a unit because lower
direct image along the raw identity functor is the identity on finite-flat
weights.
\end{proof}

\begin{theorem}[Finite-flat syntactic envelope]
\label{thm:finite-flat-syntactic-envelope}
The structure
\[
    \mathcal L_D^{\mathrm{ff}}
\]
is a locally fuzzy-domain-enriched category.  Its homs are fuzzy domains, its
composition preserves finite-flat colimits in each variable, and the Yoneda
maps
\[
    Y_{\Gamma,\Delta}:
    \mathcal L_D(\Gamma,\Delta)
    \to
    \mathcal L_D^{\mathrm{ff}}(\Gamma,\Delta)
\]
are locally fully faithful.
\end{theorem}

\begin{proof}
Each hom is a fuzzy domain by construction as an ideal completion.  Composition
is finite-flat-continuous by
Lemma~\ref{lem:finite-flat-composition-syntactic-envelope}.  Associativity and
unitality follow from
Lemma~\ref{lem:associativity-units-finite-flat-envelope}.  Local full
faithfulness of the Yoneda maps is the enriched Yoneda lemma for the principal
ideal embedding.
\end{proof}

\begin{proposition}[Affine closed structure extends]
\label{prop:affine-closed-structure-extends-to-envelope}
The symmetric affine monoidal closed structure of
\[
    \mathcal L_D
\]
extends uniquely to
\[
    \mathcal L_D^{\mathrm{ff}}.
\]
\end{proposition}

\begin{proof}
The tensor on objects is still context concatenation.  On hom fuzzy domains it
is obtained by applying external product of finite-flat weights to the raw
tensor functor
\[
    \mathcal L_D(\Gamma,\Gamma')
    \otimes
    \mathcal L_D(\Delta,\Delta')
    \longrightarrow
    \mathcal L_D(\Gamma\otimes\Delta,\Gamma'\otimes\Delta')
\]
and then taking lower direct image.

The tensor unit remains the empty context, whose terminality is preserved by
finite-flat completion because the unique raw map into the empty context
extends uniquely to the ideal-completed homs.

The closed structure is the finite-flat extension of the raw enriched
bijections
\[
    \mathcal L_D(\Gamma\otimes A,B)
    \cong
    \mathcal L_D(\Gamma,A\multimap B).
\]
Since ideal completion is functorial and locally fully faithful on principal
ideals, these raw enriched isomorphisms extend to isomorphisms of hom fuzzy
domains
\[
    \mathcal L_D^{\mathrm{ff}}(\Gamma\otimes A,B)
    \cong
    \mathcal L_D^{\mathrm{ff}}(\Gamma,A\multimap B).
\]
Thus the affine symmetric monoidal closed structure extends.
\end{proof}

\begin{proposition}[The completed syntactic reflexive object]
\label{prop:completed-syntactic-reflexive-object}
The distinguished object \(D\) of
\[
    \mathcal L_D^{\mathrm{ff}}
\]
is a linear reflexive object:
\[
    [D,D]_{\mathcal L_D^{\mathrm{ff}}}\triangleleft D.
\]
\end{proposition}

\begin{proof}
By Proposition~\ref{prop:affine-closed-structure-extends-to-envelope},
\[
    [D,D]_{\mathcal L_D^{\mathrm{ff}}}
    \cong
    D\multimap D.
\]
The raw encode and decode maps embed under Yoneda as principal finite-flat
weights, giving maps
\[
    \operatorname{encode}:D\multimap D\to D,
    \qquad
    \operatorname{decode}:D\to D\multimap D
\]
in the completed support category.  The invertible unit, the counit comparison,
and the triangle equations are the images under Yoneda of the corresponding
raw syntactic comparison cells.  Local full faithfulness of Yoneda preserves
these equations.  Hence the completed distinguished object is linearly
reflexive.
\end{proof}

\section{Models and interpretation}
\label{sec:models-and-interpretation-affine-linear-cosmos}

\begin{definition}[Linear reflexive fuzzy-domain model]
\label{def:linear-reflexive-fuzzy-domain-model-cosmos}
A \emph{linear reflexive fuzzy-domain model} is a linear reflexive fuzzy domain
\[
    M
\]
in the sense of Definition~\ref{def:linear-reflexive-fuzzy-domain-type-theory}.
\end{definition}

Such a model determines an interpretation of the syntactic envelope
\[
    \llbracket-\rrbracket_M:
    \mathcal L_D^{\mathrm{ff}}
    \longrightarrow
    \mathbf{FDom}_{\mathbb V}^{\mathbb V}.
\]
On types,
\[
    \llbracket \mathbf 1\rrbracket_M
    =
    I_{\mathrm{ff}},
\]
\[
    \llbracket D\rrbracket_M
    =
    M,
\]
\[
    \llbracket A\otimes B\rrbracket_M
    =
    \llbracket A\rrbracket_M
    \boxtimes_{\mathrm{ff}}
    \llbracket B\rrbracket_M,
\]
and
\[
    \llbracket A\multimap B\rrbracket_M
    =
    [\llbracket A\rrbracket_M,
     \llbracket B\rrbracket_M]_{\mathrm{ff}}.
\]
On contexts,
\[
    \llbracket \varnothing\rrbracket_M
    =
    I_{\mathrm{ff}},
\]
and
\[
    \llbracket x_1:A_1,\dots,x_n:A_n\rrbracket_M
    =
    \llbracket A_1\rrbracket_M
    \boxtimes_{\mathrm{ff}}
    \cdots
    \boxtimes_{\mathrm{ff}}
    \llbracket A_n\rrbracket_M.
\]

A term judgement
\[
    \Gamma\vdash t:A
\]
is interpreted as a fuzzy-domain morphism
\[
    \llbracket t\rrbracket_M:
    \llbracket \Gamma\rrbracket_M
    \longrightarrow
    \llbracket A\rrbracket_M.
\]

The primitive encode and decode terms are interpreted by the structure maps
of \(M\):
\[
    \operatorname{encode}_M:
    [M,M]_{\mathrm{ff}}\to M,
\]
\[
    \operatorname{decode}_M:
    M\to [M,M]_{\mathrm{ff}}.
\]
The derived self-application is interpreted by
\[
    \mathsf{app}_M
    =
    \operatorname{ev}_{M,M}
    \circ
    (\operatorname{decode}_M\boxtimes 1_M).
\]

\begin{theorem}[Interpretation theorem]
\label{thm:interpretation-theorem-affine-linear}
Every linear reflexive fuzzy-domain model \(M\) induces a finite-flat-continuous
affine symmetric monoidal closed locally \(\mathbb V\)-functor
\[
    \llbracket-\rrbracket_M:
    \mathcal L_D^{\mathrm{ff}}
    \longrightarrow
    \mathbf{FDom}_{\mathbb V}^{\mathbb V}
\]
preserving the distinguished linear reflexive object.
\end{theorem}

\begin{proof}
The raw term interpretation is defined by induction on terms.  Variables are
interpreted by projections obtained from weakening all unused tensor factors.
The unit term is interpreted by the unique map to \(I_{\mathrm{ff}}\).  Tensor
introduction is interpreted by the tensor product of morphisms.  Tensor
elimination is interpreted by the universal property of the tensor.  Linear
abstraction is interpreted by currying, and linear application by evaluation.
Encode and decode are interpreted by the reflexive structure maps of \(M\).

The beta, eta, reflexive beta, reflexive eta, and triangle equations are
satisfied because \(M\) is a linear reflexive fuzzy domain and
\(\mathbf{FDom}_{\mathbb V}\) is symmetric affine monoidal closed.

Thus the raw interpretation is a structure-preserving locally
\(\mathbb V\)-functor
\[
    \mathcal L_D
    \to
    \mathbf{FDom}_{\mathbb V}^{\mathbb V}.
\]
Since the homs of \(\mathcal L_D^{\mathrm{ff}}\) are finite-flat completions
of the raw homs, and the target homs \([X,Y]_{\mathrm{ff}}\) are fuzzy
domains, the universal property of \(\operatorname{Idl}_{\mathbb V}\) extends
the raw interpretation uniquely to a finite-flat-continuous interpretation on
\[
    \mathcal L_D^{\mathrm{ff}}.
\]
The extension preserves tensor, unit, internal hom, encode, and decode because
all of these structures are defined by finite-flat extension of their raw
generators.
\end{proof}

\section{Cosmos-level soundness}
\label{sec:cosmos-level-soundness}

\begin{definition}[Semantic validity of evidence judgements]
\label{def:semantic-validity-evidence-judgements}
Let \(M\) be a linear reflexive fuzzy-domain model.  A judgement
\[
    \Gamma\vdash_R s\preceq t:A
\]
is \emph{valid in \(M\)} if there is a morphism
\[
    R
    \longrightarrow
    [\llbracket \Gamma\rrbracket_M,
     \llbracket A\rrbracket_M]_{\mathrm{ff}}
    (\llbracket s\rrbracket_M,\llbracket t\rrbracket_M)
\]
in \(\mathbb V\).
\end{definition}

\begin{theorem}[Cosmos-level soundness]
\label{thm:cosmos-level-soundness}
If
\[
    \Gamma\vdash_R s\preceq t:A
\]
is derivable, then it is valid in every linear reflexive fuzzy-domain model
\(M\).
\end{theorem}

\begin{proof}
A derivation of
\[
    \Gamma\vdash_R s\preceq t:A
\]
is, by construction of the term \(\mathbb V\)-category, a morphism
\[
    R
    \longrightarrow
    \mathsf{Tm}_D(\Gamma,A)(s,t).
\]
Equivalently, it is a morphism into the relevant raw syntactic hom-object.

Applying the interpretation functor of
Theorem~\ref{thm:interpretation-theorem-affine-linear} gives a morphism
\[
    R
    \longrightarrow
    [\llbracket \Gamma\rrbracket_M,
     \llbracket A\rrbracket_M]_{\mathrm{ff}}
    (\llbracket s\rrbracket_M,\llbracket t\rrbracket_M).
\]
This is exactly semantic validity in \(M\).
\end{proof}

\section{The enriched truth lemma}
\label{sec:enriched-truth-lemma}

Let
\[
    Y:
    \mathcal L_D(\Gamma,A)
    \longrightarrow
    \mathcal L_D^{\mathrm{ff}}(\Gamma,A)
\]
be the Yoneda embedding of the raw hom \(\mathbb V\)-category into its
finite-flat completion.

For raw terms
\[
    \Gamma\vdash s,t:A,
\]
write
\[
    Ys,Yt
\]
for their images as principal finite-flat weights.

\begin{theorem}[Enriched truth lemma]
\label{thm:enriched-truth-lemma}
For every evidence object \(R\in\mathbb V\), there is a natural bijection
\[
    \mathbb V
    \bigl(
        R,
        \mathcal L_D^{\mathrm{ff}}(\Gamma,A)(Ys,Yt)
    \bigr)
    \cong
    \left\{
        \Gamma\vdash_R s\preceq t:A
    \right\}.
\]
Equivalently, the hom-object between principal syntactic terms in the
finite-flat envelope is canonically
\[
    \mathcal L_D^{\mathrm{ff}}(\Gamma,A)(Ys,Yt)
    \cong
    \mathcal L_D(\Gamma,A)(s,t).
\]
\end{theorem}

\begin{proof}
The second displayed isomorphism is the enriched Yoneda full-faithfulness of
the principal ideal embedding
\[
    Y:
    \mathcal L_D(\Gamma,A)
    \to
    \operatorname{Idl}_{\mathbb V}(\mathcal L_D(\Gamma,A)).
\]
Applying the ordinary hom-functor
\[
    \mathbb V(R,-)
\]
gives
\[
    \mathbb V
    \bigl(
        R,
        \mathcal L_D^{\mathrm{ff}}(\Gamma,A)(Ys,Yt)
    \bigr)
    \cong
    \mathbb V
    \bigl(
        R,
        \mathcal L_D(\Gamma,A)(s,t)
    \bigr).
\]
By definition of the raw term \(\mathbb V\)-category, a morphism
\[
    R\to \mathcal L_D(\Gamma,A)(s,t)
\]
is precisely a derivation of
\[
    \Gamma\vdash_R s\preceq t:A.
\]
This proves the claim.
\end{proof}

\section{Initiality and completeness}
\label{sec:initiality-completeness-cosmos}

\begin{definition}[Finite-flat affine reflexive target]
\label{def:finite-flat-affine-reflexive-target}
A \emph{finite-flat affine reflexive target} is a locally
\(\mathbb V\)-categorical category \(\mathcal E\) such that:

\begin{enumerate}
    \item every hom \(\mathbb V\)-category of \(\mathcal E\) is a fuzzy domain;

    \item composition preserves finite-flat colimits in each variable;

    \item \(\mathcal E\) is symmetric affine monoidal closed;

    \item \(\mathcal E\) contains a specified linear reflexive object
    \[
        R_{\mathcal E}
    \]
    with
    \[
        [R_{\mathcal E},R_{\mathcal E}]\triangleleft R_{\mathcal E}.
    \]
\end{enumerate}
\end{definition}

\begin{definition}[Morphisms of finite-flat affine reflexive targets]
\label{def:morphisms-finite-flat-affine-reflexive-targets}
A morphism
\[
    F:\mathcal E\to\mathcal F
\]
of finite-flat affine reflexive targets is a locally \(\mathbb V\)-functor
which:

\begin{enumerate}
    \item preserves finite-flat colimits on hom fuzzy domains;

    \item preserves the symmetric affine monoidal closed structure;

    \item sends the distinguished reflexive object of \(\mathcal E\) to the
    distinguished reflexive object of \(\mathcal F\);

    \item preserves the encode, decode, unit, counit, and triangle data of the
    reflexive projection pair.
\end{enumerate}
\end{definition}

\begin{theorem}[Initiality of the finite-flat syntactic envelope]
\label{thm:initiality-finite-flat-syntactic-envelope}
The finite-flat syntactic envelope
\[
    \mathcal L_D^{\mathrm{ff}}
\]
is the initial finite-flat affine reflexive target.

Equivalently, for every finite-flat affine reflexive target \(\mathcal E\),
there is a unique structure-preserving locally \(\mathbb V\)-functor
\[
    \mathcal L_D^{\mathrm{ff}}
    \longrightarrow
    \mathcal E.
\]
\end{theorem}

\begin{proof}
The raw syntactic category \(\mathcal L_D\) is freely generated by the affine
linear type formers, the distinguished object \(D\), the encode and decode
maps, and the specified comparison cells expressing beta, eta, reflexive beta,
reflexive eta, and the projection-pair triangle identities.

Therefore any finite-flat affine reflexive target \(\mathcal E\) receives a
unique raw structure-preserving locally \(\mathbb V\)-functor from
\[
    \mathcal L_D
\]
by interpreting the generators in \(\mathcal E\).

Since the homs of \(\mathcal E\) are fuzzy domains and composition in
\(\mathcal E\) preserves finite-flat colimits in each variable, this raw
functor extends uniquely along the local ideal completions
\[
    \mathcal L_D(\Gamma,\Delta)
    \to
    \operatorname{Idl}_{\mathbb V}(\mathcal L_D(\Gamma,\Delta))
    =
    \mathcal L_D^{\mathrm{ff}}(\Gamma,\Delta).
\]
The extension preserves composition because composition in
\(\mathcal L_D^{\mathrm{ff}}\) is defined by external product of finite-flat
weights followed by direct image, and the target composition preserves the same
finite-flat colimits.  It preserves tensor, unit, internal hom, and the
reflexive projection-pair data because these structures in
\(\mathcal L_D^{\mathrm{ff}}\) are the finite-flat extensions of the raw
generators.

Uniqueness follows from local density of principal ideals: every object of
each hom fuzzy domain in \(\mathcal L_D^{\mathrm{ff}}\) is a finite-flat
colimit of principal syntactic substitutions.
\end{proof}

\begin{corollary}[Cosmos-level completeness]
\label{cor:cosmos-level-completeness}
Let
\[
    \Gamma\vdash s,t:A
\]
be raw terms.  If every linear reflexive fuzzy-domain model validates an
evidence comparison from \(R\) between \(s\) and \(t\), then the classifying
syntax contains the corresponding comparison morphism
\[
    R
    \longrightarrow
    \mathcal L_D^{\mathrm{ff}}(\Gamma,A)(Ys,Yt).
\]
For principal syntactic terms, this comparison is represented uniquely by a
derivation
\[
    \Gamma\vdash_R s\preceq t:A.
\]
\end{corollary}

\begin{proof}
Semantic validity in all models includes validity in the initial model
\[
    \mathcal L_D^{\mathrm{ff}}.
\]
Thus there is a morphism
\[
    R
    \to
    \mathcal L_D^{\mathrm{ff}}(\Gamma,A)(Ys,Yt).
\]
By the enriched truth lemma, this is equivalently a derivation
\[
    \Gamma\vdash_R s\preceq t:A.
\]
\end{proof}

\section{Morphism of concrete reflexive fuzzy-domain models}
\label{sec:morphisms-concrete-reflexive-models}

For concrete models, the correct notion of morphism uses projection pairs,
because a single morphism
\[
    D\to E
\]
does not induce a canonical map
\[
    [D,D]_{\mathrm{ff}}\to [E,E]_{\mathrm{ff}}.
\]
A projection pair does.

\begin{definition}[Induced internal-hom projection pair]
\label{def:induced-internal-hom-projection-pair}
Let
\[
    (e,p):D\triangleleft E
\]
be a fuzzy projection pair.  Define
\[
    \widehat e:[D,D]_{\mathrm{ff}}\to [E,E]_{\mathrm{ff}}
\]
by
\[
    \widehat e(h)=e\,h\,p,
\]
and define
\[
    \widehat p:[E,E]_{\mathrm{ff}}\to [D,D]_{\mathrm{ff}}
\]
by
\[
    \widehat p(k)=p\,k\,e.
\]
Then
\[
    (\widehat e,\widehat p):
    [D,D]_{\mathrm{ff}}
    \triangleleft
    [E,E]_{\mathrm{ff}}
\]
is the induced internal-hom projection pair.
\end{definition}

\begin{proof}
The maps preserve finite-flat colimits by separate finite-flat-continuity of
composition.  Moreover,
\[
    \widehat p\widehat e(h)
    =
    p(ehp)e
    \cong
    h
\]
using the invertible unit
\[
    pe\cong \id_D.
\]
For
\[
    k:E\to E,
\]
one has
\[
    \widehat e\widehat p(k)
    =
    e(pke)p
    \cong
    (ep)k(ep)
    \Rightarrow
    k
\]
using the counit comparison
\[
    ep\Rightarrow\id_E.
\]
The triangle identities follow from those of \((e,p)\) and the functoriality
of composition.
\end{proof}

\begin{definition}[Projection-pair morphism of linear reflexive models]
\label{def:projection-pair-morphism-linear-reflexive-models}
Let
\[
    \mathcal M=(D,\operatorname{encode}_D,\operatorname{decode}_D)
\]
and
\[
    \mathcal N=(E,\operatorname{encode}_E,\operatorname{decode}_E)
\]
be linear reflexive fuzzy-domain models.  A morphism
\[
    \mathcal M\to\mathcal N
\]
is a fuzzy projection pair
\[
    (e,p):D\triangleleft E
\]
such that, with
\[
    (\widehat e,\widehat p):
    [D,D]_{\mathrm{ff}}\triangleleft [E,E]_{\mathrm{ff}}
\]
as in Definition~\ref{def:induced-internal-hom-projection-pair}, there are
coherent comparison isomorphisms
\[
    e\,\operatorname{encode}_D
    \cong
    \operatorname{encode}_E\,\widehat e,
\]
and
\[
    \widehat p\,\operatorname{decode}_E
    \cong
    \operatorname{decode}_D\,p.
\]
The comparison cells are required to be compatible with the unit, counit, and
triangle identities of the reflexive projection pairs.
\end{definition}

\section{Extensional reflection}
\label{sec:extensional-reflection-cosmos}

\begin{definition}[Extensional syntax]
\label{def:extensional-syntax-cosmos}
The \emph{extensional affine reflexive syntax} is obtained from the raw syntax
by adjoining the reverse reflexive eta comparison
\[
    d
    \preceq
    \operatorname{encode}(\operatorname{decode}(d))
    :
    D.
\]
Equivalently, it adjoins
\[
    d
    \preceq
    \lambda^D x.(d\bullet x)
    :
    D.
\]
At the cosmos level, this is the addition of a global comparison cell
\[
    I
    \longrightarrow
    \mathsf{Tm}_D(\Gamma,D)
    \bigl(d,\operatorname{encode}(\operatorname{decode}(d))\bigr)
\]
natural in \(\Gamma\) and \(d\).
\end{definition}

Let
\[
    \mathcal L_D^{\mathrm{ext}}
\]
be the raw locally \(\mathbb V\)-categorical syntax with this additional
comparison cell and its coherence equations.  Let
\[
    (\mathcal L_D^{\mathrm{ext}})^{\mathrm{ff}}
\]
be its finite-flat syntactic envelope.

\begin{theorem}[Extensional initiality]
\label{thm:extensional-initiality-cosmos}
The finite-flat envelope
\[
    (\mathcal L_D^{\mathrm{ext}})^{\mathrm{ff}}
\]
is the initial finite-flat affine reflexive target whose distinguished
reflexive object is extensional.
\end{theorem}

\begin{proof}
The proof is the same as
Theorem~\ref{thm:initiality-finite-flat-syntactic-envelope}, using the
extended enriched presentation.  In any extensional target, the semantic
counit
\[
    \operatorname{encode}\operatorname{decode}
    \Rightarrow
    \id_D
\]
is invertible, so the added reverse comparison is interpreted by its inverse.
Conversely, the added comparison makes the syntactic counit invertible in the
initial target.
\end{proof}

\begin{corollary}[Extensional soundness and completeness]
\label{cor:extensional-soundness-completeness-cosmos}
For extensional linear reflexive fuzzy-domain models, the extensional syntax is
sound and complete in the sense of
Theorem~\ref{thm:cosmos-level-soundness} and
Corollary~\ref{cor:cosmos-level-completeness}.
\end{corollary}

\section{The posetal reflection}
\label{sec:posetal-reflection-affine-linear-type-theory}

Let
\[
    \Lambda=(L,\leq,\wedge,\vee,\otimes,\multimap,\bot,\top)
\]
be a fuzzy truth-value object, regarded as a posetal fuzzy cosmos.  In this
case an evidence object is a truth value
\[
    a\in L,
\]
and a morphism
\[
    a\to C(x,y)
\]
is precisely the inequality
\[
    a\leq C(x,y).
\]

Thus the cosmos-indexed judgement
\[
    \Gamma\vdash_R s\preceq t:A
\]
specializes to the scalar fuzzy judgement
\[
    \Gamma\vdash_a s\preceq t:A.
\]

The structural rules become:

\[
    \frac{}
         {\Gamma\vdash_{\top} t\preceq t:A},
\]

\[
    \frac{\Gamma\vdash_a r\preceq s:A
          \qquad
          \Gamma\vdash_b s\preceq t:A}
         {\Gamma\vdash_{b\otimes a} r\preceq t:A},
\]

\[
    \frac{\Gamma\vdash_b s\preceq t:A
          \qquad
          a\leq b}
         {\Gamma\vdash_a s\preceq t:A},
\]

and
\[
    \frac{\Gamma\vdash_{a_i} s\preceq t:A\quad(i\in I)}
         {\Gamma\vdash_{\bigvee_{i\in I}a_i} s\preceq t:A}.
\]

The congruence rules specialize to the expected tensor-degree rules:
\[
    \frac{\Gamma\vdash_a s\preceq s':A
          \qquad
          \Delta\vdash_b t\preceq t':B}
         {\Gamma,\Delta\vdash_{a\otimes b}
          s\otimes t\preceq s'\otimes t':A\otimes B},
\]
\[
    \frac{\Gamma\vdash_a f\preceq g:A\multimap B
          \qquad
          \Delta\vdash_b u\preceq v:A}
         {\Gamma,\Delta\vdash_{a\otimes b}
          f\,u\preceq g\,v:B},
\]
and similarly for encode, decode, abstraction, and self-application.

\begin{definition}[Posetal semantic validity]
\label{def:posetal-semantic-validity-affine-linear}
For a linear reflexive \(\Lambda\)-domain model \(M\), write
\[
    M\models \Gamma\vdash_a s\preceq t:A
\]
if
\[
    a
    \leq
    [\llbracket \Gamma\rrbracket_M,
     \llbracket A\rrbracket_M]_{\mathrm{ff}}
    (\llbracket s\rrbracket_M,\llbracket t\rrbracket_M).
\]
Write
\[
    \Gamma\models_a s\preceq t:A
\]
if this holds in every linear reflexive \(\Lambda\)-domain model.
\end{definition}

\begin{theorem}[Posetal soundness and completeness]
\label{thm:posetal-soundness-completeness-affine-linear}
For every context \(\Gamma\), terms
\[
    \Gamma\vdash s,t:A,
\]
and truth value \(a\in L\),
\[
    \Gamma\vdash_a s\preceq t:A
    \quad\Longleftrightarrow\quad
    \Gamma\models_a s\preceq t:A.
\]
\end{theorem}

\begin{proof}
The forward implication is the posetal specialization of
Theorem~\ref{thm:cosmos-level-soundness}.

For the reverse implication, evaluate the judgement in the initial
finite-flat syntactic envelope.  By the enriched truth lemma,
\[
    a
    \leq
    \mathcal L_D^{\mathrm{ff}}(\Gamma,A)(Ys,Yt)
\]
holds if and only if there is a derivation
\[
    \Gamma\vdash_a s\preceq t:A.
\]
Thus semantic validity in all models implies derivability.
\end{proof}

\begin{definition}[Top-degree definitional equality]
\label{def:top-degree-definitional-equality-posetal}
In the posetal reflection, define
\[
    \Gamma\vdash s\equiv t:A
\]
to mean that both
\[
    \Gamma\vdash_{\top} s\preceq t:A
\]
and
\[
    \Gamma\vdash_{\top} t\preceq s:A
\]
are derivable.
\end{definition}

Ordinary beta and eta for the genuine linear function type are definitional
equalities.  Reflexive beta is a definitional equality.  Reflexive eta is
generally only the one-sided approximation
\[
    \Gamma
    \vdash_{\top}
    \operatorname{encode}(\operatorname{decode}(d))
    \preceq
    d
    :
    D.
\]
In extensional models, the reverse approximation is also valid and derivable
in the extensional reflection.

\section{The crisp cartesian specialization}
\label{sec:crisp-cartesian-specialization-affine-linear}

Let
\[
    \Lambda=2.
\]
In the separated case, fuzzy domains specialize to dcpos, and fuzzy-domain
morphisms specialize to Scott-continuous maps.  Finite-flat weights specialize
to ordinary directed ideals.

In this crisp separated specialization, the tensor product of fuzzy domains
specializes to the cartesian product of dcpos:
\[
    A\boxtimes_{\mathrm{ff}}B
    =
    A\times B.
\]
Thus every object has a diagonal
\[
    X\longrightarrow X\times X.
\]
Consequently contraction becomes sound after passing to the cartesian
specialization.

A linear reflexive fuzzy domain becomes an ordinary reflexive dcpo
projection-pair situation
\[
    [D,D]\triangleleft D.
\]
If the projection pair is extensional, then
\[
    D\cong [D,D].
\]
The derived operations
\[
    s\bullet t
    =
    \operatorname{decode}(s)(t)
\]
and
\[
    \lambda^D x.t
    =
    \operatorname{encode}(\lambda x.t)
\]
recover the usual reflexive-domain interpretation of the untyped lambda
calculus.

The general fuzzy-cosmos theory is therefore not a notational deformation of
cartesian lambda calculus.  Its affine-linear discipline is forced by the
absence of canonical contraction maps in the symmetric affine monoidal closed
category of fuzzy domains.

\section{Summary}
\label{sec:summary-affine-linear-type-theory}

This chapter constructed the affine linear type theory associated to a
reflexive fuzzy domain in the fuzzy cosmos setting.

The semantic base is the locally \(\mathbb V\)-categorical category of fuzzy
domains
\[
    \mathbf{FDom}_{\mathbb V}^{\mathbb V},
\]
with hom fuzzy domains
\[
    [A,B]_{\mathrm{ff}}.
\]
Its tensor is affine, not cartesian.  Hence weakening is available, but
contraction is not.

The primitive proof-theoretic comparison judgement is cosmos-indexed:
\[
    \Gamma\vdash_R s\preceq t:A,
    \qquad
    R\in\mathbb V.
\]
It is interpreted as a morphism
\[
    R
    \longrightarrow
    [\llbracket \Gamma\rrbracket,
     \llbracket A\rrbracket]_{\mathrm{ff}}
    (\llbracket s\rrbracket,\llbracket t\rrbracket)
\]
in the fuzzy cosmos.

The raw syntax forms a locally \(\mathbb V\)-categorical symmetric affine
monoidal closed category
\[
    \mathcal L_D
\]
with a distinguished linear reflexive object \(D\).  Applying the fuzzy ideal
monad locally gives the finite-flat syntactic envelope
\[
    \mathcal L_D^{\mathrm{ff}}(\Gamma,\Delta)
    =
    \operatorname{Idl}_{\mathbb V}
    \bigl(\mathcal L_D(\Gamma,\Delta)\bigr).
\]
This envelope is the initial finite-flat affine closed locally
\(\mathbb V\)-categorical category equipped with a linear reflexive object.

Soundness follows by interpreting the initial syntax in any linear reflexive
fuzzy-domain model.  Completeness follows from the enriched truth lemma:
for raw terms \(s,t\),
\[
    \mathcal L_D^{\mathrm{ff}}(\Gamma,A)(Ys,Yt)
    \cong
    \mathcal L_D(\Gamma,A)(s,t).
\]
In the posetal reflection over a fuzzy truth-value object \(\Lambda\), this
becomes the scalar theorem
\[
    \Gamma\vdash_a s\preceq t:A
    \quad\Longleftrightarrow\quad
    \Gamma\models_a s\preceq t:A.
\]

In the crisp separated specialization \(\Lambda=2\), the affine tensor becomes
the cartesian product of dcpos, contraction becomes sound, and the
distinguished \(D\)-fragment recovers the ordinary reflexive-domain semantics
of the untyped lambda calculus.